% Reunited development from DESARROLLO_MATEMATICO, sections 38--40 and 50.
\subsubsection{Capacity index, vacancies, and induction of the tritic regime}
\label{vac-capacidad-trit}

The vacancy sequence is linked to refinement through an increment identity.
To exhibit it, fix the phase origin \(\theta=0\). Let
\[
 \rho=\log_{10}9=1-\delta,\qquad
 \mathcal C_t=\lfloor(t+1)\rho\rfloor\quad(t\geq0).
\]
The index \(\mathcal C_t\) expresses the decimal capacity of refinement
in the normalization comparing a ternary block of \(6\) positions,
with \(729\) possibilities, and a decimal block of \(3\) positions,
with \(1000\) possibilities. Indeed,
\(\log_{1000}729=\log_{10}9=\rho\). The letter \(K\) is reserved
for the dodecaphase record constructed later.

\begin{proposition}[Complementarity of capacity and vacancy]
\label{prop:capacidad-vacancia}
For \(t\geq1\),
\begin{equation}
 \mathcal C_t-\mathcal C_{t-1}=1-z_t(0).
 \label{eq:capacidad-vacancia}
\end{equation}
Consequently, if \(0\leq a<b\),
\[
 \mathcal C_b-\mathcal C_a+
 \sum_{t=a+1}^{b}z_t(0)=b-a.
\]
\end{proposition}
\begin{proof}
Irrationality of \(\delta\) gives
\[
 \mathcal C_t
 =t+1-\lceil(t+1)\delta\rceil
 =t-\lfloor(t+1)\delta\rfloor.
\]
Subtracting consecutive terms produces \eqref{eq:capacidad-vacancia}.
Summation over the interval telescopes.
\end{proof}

Each transition increases the block counter by one. If \(z_t=0\),
capacity also increases; if \(z_t=1\), capacity remains unchanged
while a transition is incorporated into the history. Here a vacancy
is a capacity-retention event, with a determined position and antecedent.
The preceding balance preserves length exactly between acquired capacity
and recorded vacancies. The observed phase \(g_t^{\rm cap}=[\mathcal C_t]_9\)
satisfies
\[
 g_t^{\rm cap}-g_{t-1}^{\rm cap}\equiv1-z_t\pmod9.
\]
This capacity phase and chronological phase \([t]_9\) are two
different observations. Local capacity retention and a complete holonomy
turn have distinct mechanisms, although both allow the phase observation
to be preserved while the state changes.

The relation with TRIT can be followed in the vacancy positions themselves.
For \(j\geq1\), define
\[
 p_j=\left\lfloor\frac j\delta\right\rfloor,\qquad
 v_j=\mathcal C_{p_j},\qquad h_j=p_{j+1}-p_j.
\]
The inequality \(p_j\delta<j<(p_j+1)\delta\) identifies \(p_j\)
as the position of event \(j\). Since \(0<\delta<1\), it also
gives \(\lfloor(p_j+1)\delta\rfloor=j\), and hence
\[
 v_j=p_j-j,\qquad
 v_{j+1}-v_j=h_j-1.
\]
If \(s_j=22-h_j\), then \(s_j\in\{0,1\}\) and
\begin{equation}
 [v_{j+1}]_3=[v_j-s_j]_3.
 \label{eq:vacancia-regimen}
\end{equation}
An interval of length \(22\) preserves the class; one of length
\(21\) shifts it by one in the negative direction. The canonical
selector is applied to the positive phase \(r_j=1+[v_j]_9\):
\[
 \tau_j=M_{\rm ph}(r_j).
\]
A tritic regime is thus determined by the capacity index retained
at each event, in addition to the vacancy's binary encoding.

The positions of the short intervals form the induced mechanical
coding of slope \(22-1/\delta\). Their separations \(6\) and \(7\)
determine the lengths of plateaus of the class \([v_j]_3\), once
the origin and endpoint convention are fixed. The first classes are
\(2\) for \(6\) events, \(1\) for \(7\), and \(0\) for \(7\);
the first segment of \(6\) is the initial restriction of a plateau.
Their signatures traverse \(0,-1,+1\). Equations
\eqref{eq:capacidad-vacancia} and \eqref{eq:vacancia-regimen}
make explicit the link between refinement, vacancies, and regime
selection. The quadratic reading of TRIT subsequently realizes these
types through their respective generators; the temporal event and
its geometric representation thus retain an identifiable dependence.

Inverting the monotone index supplies another useful observation:
\[
 \mathcal T(k)=\min\{t:\mathcal C_t=k\}
   =\left\lceil\frac k\rho\right\rceil-1\qquad(k\geq1).
\]
If \(\sigma=\rho^{-1}-1\), the number of blocks needed to gain
\(m\) units of capacity is
\[
 \mathcal T(k+m)-\mathcal T(k)
 =m+\lceil(k+m)\sigma\rceil-\lceil k\sigma\rceil
 \in\{m+\lfloor m\sigma\rfloor,m+\lceil m\sigma\rceil\}.
\]
In particular, \(9\) units of capacity require \(9\) or \(10\)
blocks. The period of a finite phase remains exact, whereas return
time measured by the other index admits two values. This distinction
belongs to refinement dynamics and is not an a posteriori correction
of its numerical values.
